\chapter{Objects, values, and native data structures}
\label{ch:python-data}

\begin{learningobjectives}
After this chapter, you should be able to:
\begin{itemize}
  \item distinguish names, objects, types, values, identity, and mutability;
  \item choose strings, sequences, mappings, and sets from their semantics;
  \item preserve raw data while creating validated representations; and
  \item use iteration without confusing traversal with materialization.
\end{itemize}
\end{learningobjectives}

\begin{chapterconnection}
Chapter 2 established which Python process is running. We now look inside that
process. The central question changes from ``Which interpreter owns this
computation?'' to ``What object does this name refer to, and which operations
preserve its meaning?''
\end{chapterconnection}

\section{Names refer to objects}

Python assignment binds a name to an object. It does not declare a fixed box
whose type can never change.

\begin{lstlisting}
probability_text = "0.873"
probability = float(probability_text)

assert isinstance(probability_text, str)
assert isinstance(probability, float)
\end{lstlisting}

The original string and parsed floating-point value are different objects with
different meanings. Preserving both supports diagnosis and provenance. A value
that looks numerical is not numerical merely because a human can interpret it.

\begin{definitionbox}{type}
A type defines a family of objects and the operations and behaviors associated
with them. A type annotation documents an intended contract; runtime validation
is a separate action unless a library or explicit check performs it.
\end{definitionbox}

Equality asks whether two objects represent equal values under their type's
rules. Identity asks whether two names refer to the same object. Use
\code{is None} for the singleton \code{None}; do not replace value equality with
identity comparisons.

\section{Mutability and aliasing}

An immutable object cannot change its own value after creation. A mutable object
can. If two names refer to one mutable list, a change through either name is
visible through both.

\begin{lstlisting}
temperatures = [20.1, 20.4]
shared_reference = temperatures
independent_copy = temperatures.copy()

shared_reference.append(20.8)

assert temperatures == [20.1, 20.4, 20.8]
assert independent_copy == [20.1, 20.4]
\end{lstlisting}

Aliasing is not automatically a bug. It is a state-sharing decision. Bugs arise
when code shares mutable state without making ownership visible.

\section{Strings are structured text}

A \term{string} is an immutable sequence of Unicode text. Text arrives from
filenames, identifiers, logs, forms, CSV fields, JSON keys, command-line
arguments, and API payloads. It may contain invisible whitespace, distinct
Unicode representations, or characters that resemble digits without having
numeric type.

Reliable text work separates stages:

\begin{center}
\begin{tikzpicture}[node distance=8mm]
  \node[wideflowbox] (raw) {Raw evidence\\preserved exactly};
  \node[wideflowbox,right=of raw] (parse) {Contract-aware\\parsing and normalization};
  \node[wideflowbox,right=of parse] (typed) {Typed value\\with explicit meaning};
  \draw[flowarrow] (raw) -- (parse);
  \draw[flowarrow] (parse) -- (typed);
\end{tikzpicture}
\end{center}

Normalization is a policy, not a cleaning reflex. Lowercasing may be correct for
a case-insensitive category and destructive for a case-sensitive identifier.
Formatting a number for display must not replace the full-precision value used
in later computation.

\section{Sequences encode order}

Lists and tuples are both ordered sequences. A list is mutable and is suitable
for a collection that grows or changes under clear ownership. A tuple is
immutable and often communicates a fixed record, coordinate, or return bundle.

Indexing retrieves one element. Slicing constructs a subsequence. Negative
indices count from the end. These operations are convenient, but meaning still
comes from the data contract. The third element of an unlabeled tuple is easy to
misinterpret in a large system.

\begin{professionalbox}
Choose a sequence because order and multiplicity are meaningful, not because a
list is the first container you remember. For domain records, a dataclass or
validated model often communicates more than positional elements.
\end{professionalbox}

\section{Mappings and sets encode different relations}

A dictionary maps unique keys to values. It is appropriate for labeled fields,
indexes, configuration, counts, and adjacency lists. A set represents unique
membership and supports union, intersection, and difference.

\begin{lstlisting}
sensor = {
    "sensor_id": "T-17",
    "unit": "degree Celsius",
    "precision": 0.1,
}

observed = {"temperature", "pressure"}
required = {"temperature", "humidity"}
missing = required - observed

assert missing == {"humidity"}
\end{lstlisting}

A missing key, a key explicitly mapped to \code{None}, and a key mapped to an
empty string are three distinct states. Conflating them erases information.

\section{Iterables, iterators, and generators}

An \term{iterable} can produce an iterator. An \term{iterator} yields one value
at a time and remembers its position. A \term{generator} is a convenient way to
implement an iterator with suspended function state.

Iteration permits bounded, streaming computation. It does not guarantee that
the source is small, repeatable, or safe to consume twice. Converting an
iterator to a list materializes every remaining item in memory.

\begin{lstlisting}
def valid_probabilities(lines):
    for line in lines:
        value = float(line)
        if 0.0 <= value <= 1.0:
            yield value
\end{lstlisting}

This generator can process a file progressively. It still needs a policy for
malformed input, provenance, logging, and whether silently excluding values is
scientifically acceptable.

\begin{warningbox}
Truthiness is not a missing-data policy. Zero, an empty string, an empty list,
\code{False}, and \code{None} are all false in Boolean context but may represent
very different scientific states.
\end{warningbox}

\begin{checkpointbox}
Design a representation for one experiment with an identifier, ordered time
points, unique observed variables, and metadata fields. Justify each container
from its semantics and identify which state is mutable.
\end{checkpointbox}

\section{Worked example: from instrument text to a trustworthy record}

Suppose an instrument exports this line:

\begin{lstlisting}
raw_line = "  SAMPLE-017 , 23.40 , valid  "
\end{lstlisting}

Humans see an identifier, temperature, and quality label. Python initially sees
one string. We should not jump directly to a cleaned list and discard the
original. Each transformation answers a separate question.

\begin{workedexample}{build meaning in observable stages}
\textbf{Stage 1: preserve evidence.}

\begin{lstlisting}
raw_line = "  SAMPLE-017 , 23.40 , valid  "
fields = raw_line.split(",")
assert fields == ["  SAMPLE-017 ", " 23.40 ", " valid  "]
\end{lstlisting}

Splitting is a syntactic operation. It assumes comma is the delimiter but has
not yet decided what any field means.

\textbf{Stage 2: normalize under a field-specific contract.}

\begin{lstlisting}
sample_id_text, temperature_text, quality_text = fields
sample_id = sample_id_text.strip()
quality = quality_text.strip().casefold()
\end{lstlisting}

Whitespace is removed because the export contract declares it insignificant.
The quality label is case-insensitive. We do not case-normalize the identifier
without an equivalent promise.

\textbf{Stage 3: convert and validate.}

\begin{lstlisting}
temperature_c = float(temperature_text)

if not -273.15 <= temperature_c <= 2_000.0:
    raise ValueError("temperature_c is outside the accepted domain")
if quality not in {"valid", "suspect", "failed"}:
    raise ValueError("quality has an unknown label")
\end{lstlisting}

Conversion answers ``can these characters represent a floating-point value?''
The range check answers a different domain question. The upper limit here is an
example contract, not a law Python can infer.

\textbf{Stage 4: construct a labeled record.}

\begin{lstlisting}
record = {
    "sample_id": sample_id,
    "temperature_c": temperature_c,
    "quality": quality,
    "raw_line": raw_line,
}
\end{lstlisting}

The dictionary makes field roles visible and retains the evidence needed to
diagnose parsing. A later chapter will replace suitable dictionaries with
validated domain classes.
\end{workedexample}

Notice what the code does not establish: whether the sensor was calibrated,
whether the sample identifier belongs to the experiment, or whether
\code{23.40} was truly reported in degrees Celsius. Representation checks and
scientific validity remain distinct.

\section{Choosing a container by the question it answers}

Container choice is clearer when phrased as a question:

\begin{tabularx}{\textwidth}{@{}p{0.25\textwidth}X@{}}
\toprule
Structure & Question it answers naturally \\
\midrule
String & What text, identifier, or encoded field was supplied? \\
List & What ordered collection may be updated under one owner? \\
Tuple & What fixed ordered group should not change? \\
Dictionary & Which value belongs to this unique key or field name? \\
Set & Is this distinct item present, and how do memberships compare? \\
Iterator & What is the next item without materializing the entire source? \\
\bottomrule
\end{tabularx}

No table can choose without domain context. A coordinate might be a tuple; a
mutable simulation state might need a class; an observation table will later
become a DataFrame. The principle is to make important semantics visible.

\begin{misconceptionbox}
\textbf{Misconception:} tuples are ``safer lists'' and should replace lists
whenever mutation is not expected.

\textbf{Correction:} immutability is useful, but positional meaning can remain
opaque. A tuple of three floats might be a point, RGB color, confidence
interval, or date. When field names and invariants matter, a dataclass or
validated model may communicate the domain better.
\end{misconceptionbox}

\section{Iteration changes when work happens}

Compare a list comprehension and generator expression:

\begin{lstlisting}
squares_list = [value**2 for value in range(1_000_000)]
squares_stream = (value**2 for value in range(1_000_000))
\end{lstlisting}

The first computes and stores one million results immediately. The second
creates an iterator that computes results as requested. That difference affects
memory, timing, errors, and repeatability. A generator may read from a file that
changes, raise an exception halfway through, or be exhausted after one pass.

``Lazy'' does not mean ``free.'' It means computation is deferred and therefore
must be owned by the consumer that triggers it.

\conceptcheck{If \code{values} is a generator, why can \code{sum(values)} work
once while a second \code{sum(values)} returns zero? What observation would
confirm your explanation?}

\section{Exercises}

\subsection*{Foundational}
\begin{enumerate}
  \item Explain the difference between \code{==} and \code{is}. Give one correct
  use of \code{is} from this chapter.
  \item Give distinct scientific meanings for \code{0}, \code{""},
  \code{False}, and \code{None}. Explain why one truthiness check cannot replace
  a missing-data policy.
  \item For each native container, state whether order, duplicates, labels, and
  mutation are central to its contract.
\end{enumerate}

\subsection*{Applied}
\begin{enumerate}
  \item Extend the instrument example with a timestamp field. Preserve raw text,
  parse a typed value, and name two validity questions parsing cannot answer.
  \item Given repeated experiment tags, use a set to report distinct tags and a
  dictionary to count them. Explain why the two results answer different
  questions.
  \item Write a generator that yields only nonblank lines with their original
  line numbers. Decide whether whitespace-only lines are missing data or
  ignorable formatting, and document that policy.
\end{enumerate}

\subsection*{Design}
Choose a scientific object currently represented as a list or dictionary in
your own work. Defend that representation or propose a clearer one using order,
identity, mutability, labels, and invariants as criteria.

\begin{chapterreview}
\textbf{Key ideas.} Names refer to objects. Types determine behavior. Mutable
state requires visible ownership. Text normalization is field-specific policy.
Containers encode order, labels, uniqueness, or traversal. Parsing builds a
representation but cannot establish scientific truth.

\textbf{Vocabulary.} object; name; type; value; identity; equality; mutability;
aliasing; Unicode; sequence; mapping; set; iterable; iterator; generator.

\textbf{Explain aloud.} Why is preserving \code{raw\_line} useful even after a
record has been successfully parsed?
\end{chapterreview}

\begin{notebooklink}
Lecture 1 notebooks 01 through 05 develop these ideas through detailed examples
on objects, strings, sequences, dictionaries, sets, comprehensions, and
generators.
\end{notebooklink}

\section{Takeaway}

Python's native data structures are modeling decisions. Reliable work begins by
matching representation to meaning, preserving evidence, and making state
ownership explicit.
